\documentclass[11pt]{letter}
\usepackage[T1]{fontenc}
\usepackage{newtxtext}
\usepackage[margin=0.85in]{geometry}

\signature{Yongzhe Li\\Corresponding author\\School of Mechanical Engineering\\Southeast University}
\address{School of Mechanical Engineering\\Southeast University\\Nanjing 211189, China\\yongzheli@seu.edu.cn}

\begin{document}

\begin{letter}{Editor-in-Chief\\Expert Systems with Applications}
\opening{Dear Editor,}

Please consider the manuscript, ``Reinforcement learning with forecast-loss rewards for
sensor scheduling under power constraints,'' for publication as a
Research Article in \emph{Expert Systems with Applications}.

The paper presents the Prediction-Driven Proximal Policy Optimization (PD-PPO)
scheduler for a power budget that prevents all channels from operating
simultaneously. Its complete training configuration combines a
reward based on forecast loss, behavior cloning pretraining, continuing
advantage-weighted behavior cloning, an auxiliary future-condition
classification loss, and
simulated warning signals. A feasibility layer checks the steady-state and
startup-peak power limits, while the minimum dwell-time constraint
determines the executed channel subset. All six candidates satisfy both power
limits in the reported $q=1$ instance, so the optional-channel cardinality
constraint and minimum dwell-time constraint are the active restrictions.

Across 22 post-selection evaluation seeds in a controlled six-channel
cold-region monitoring simulation, the complete PD-PPO configuration achieved
10.6\% lower mean forecast loss and 8.0\% lower macro-averaged normalized
forecast loss than the validation-selected fixed-schedule baseline. Both losses are
lower in every seed and remain lower over each continuous 5,242-step scoreable
test partition. The separate 24-seed descriptive aggregate, which includes two
pilot/model-selection seeds, retains the 10.1\% and 7.9\% reductions. The
evaluation also includes matched PPO training configurations with different
scalar objectives, a Double DQN training-configuration comparison, an
alternative-forecaster analysis, and executed-action diagnostics. The
privileged one-step look-ahead reference uses next-step test targets. The
privileged true-label baseline uses the exact label sequence
$c_t,\ldots,c_{t+16}$ at each decision; neither information source is available
to a deployed scheduler.

The manuscript fits the journal's scope through constrained sensor scheduling,
forecast-based evaluation, explicit information boundaries, and reproducible
system-level analysis. It is original, is not under consideration elsewhere,
and has been approved by both authors. The authors declare no competing interests.

The code and evidence repository will be made public upon acceptance. No
graphical asset was created or modified using a generative image tool. The
anonymous manuscript, title page, highlights, supplementary material, and
reproducibility archive are separate submission files.

\closing{Sincerely,}

\end{letter}
\end{document}
